\cotefile{29}
\scanpage{49}
\mspage{- 23 -}

n$^{\circ}$385

\uline{l'injection de} $Y$ \uline{dans} $X$. \uline{Pour tout point} $a$ \uline{de} $Y$, \uline{l'application} $\iota_a$ \uline{est une bijection de} $D_a(Y)$ \uline{sur le sous-espace de} $D_a(X)$ \uline{constitué des} $\mu$ \uline{telles que} $f.\mu=0$ \uline{pour toute fonction de classe} $C^r$ \uline{dans un voisinage} $U$ \uline{de} $a$, \uline{nulle en tout point de} $Y\cap U$.

Soit $\mu$ dans $D_a(Y)$ et $f$ de classe $C^r$ dans un voisinage $U$ de $a$, nulle sur $Y\cap U$ ; on a :
\[
 f.\iota_a(\mu)=\iota_a((f\circ\iota).\mu)=0
\]
puisque $f\circ\iota=0$. Comme $Y$ est une sous-variété de $X$, il existe un ouvert $U$ contenant $a$, un système de coordonnées $x_1,\ldots,x_n$ dans $U$, et un entier $m$ tel que $\handwritten{0}\leq m\leq n$ et que $Y\cap U$ soit l'ensemble des points de $U$ où s'annulent $x_{m+1},x_{m+2},\ldots,x_n$. Les fonctions $x_1\circ\iota,\ldots,x_m\circ\iota$ forment alors un système de coordonnées sur l'ouvert $Y\cap U$ de $Y$.

L'espace $D_a(Y)$ a une base $\{\Lambda^\alpha\}_{\alpha\in N_r^m}$ caractérisée par :
\[
 \langle\Lambda^\alpha,x^{\alpha'}\circ\iota\rangle=\delta_{\alpha\alpha'}\qquad \alpha,\alpha'\in N_r^m
\]
et l'espace $D_a(X)$ admet une base $\{\Delta^\beta\}_{\beta\in N_r^n}$ définie par :
\[
 \langle\Delta^\beta,x^{\beta'}\rangle=\delta_{\beta\beta'}\qquad \beta,\beta'\in N_r^n.
\]

Nous noterons $0$ l'élément nul de $N_r^{n-m}$ ; on a alors $\iota_a(\Lambda^\alpha)=\Delta^{(\alpha,0)}$ pour tout $\alpha$ dans $N_r^m$ et par conséquent, $\iota_a$ est injective.

Par ailleurs, soit $\mu$ dans $D_a(X)$ telle que $x_i.\mu=0$ pour $m<i\leq n$. D'après (\unsure{5|3}9), on en déduit $\langle\mu,x^\beta\rangle=0$ pour tout multiindice $\handwritten{\beta}$ dans $N_r^n$ pour lequel il existe $i$ avec $m<i\leq n$ et $\handwritten{\beta}_i\ne 0$ ; autrement dit, on ne peut avoir $\langle\mu,x^\beta\rangle\ne 0$ que si $\handwritten{\beta}=(\alpha,0)$ pour $\alpha$ dans $N_r^m$ convenable. De la formule $\mu=\sum\limits_\beta\langle\mu,x^\beta\rangle.\Delta^\beta$, on déduit alors $\mu=\sum\limits_\alpha\langle\mu,x^\alpha\rangle\,\iota_a(\Lambda^\alpha)$, ce qui prouve que $\mu$ est dans l'image de $\iota_a$.

\hfill C.Q.F.D.

Pour tout point $a$ de $X$, l'espace vectoriel $D_a(X)\otimes D_a(X)$ est muni d'une structure
% archive sheet number: 48
\scanpage{50}
\mspage{II}

Soit $Y_1$ le normalisé de $Y$ en $y$, \inserted{\unsure{pt dés. point mt|pt dés. point} \unsure{chlt|clôt.}} \deleted{\unsure{Alors}}, en \unsure{le pt|tt} $y\in Y$. \unsure{Quitte} à se \deleted{localiser} \unsure{en} un pt maximal de la fibre \deleted{\unsure{spéciale}} \inserted{de $X_y$}, on \deleted{\unsure{prouve}} peut supposer $X_y$ \deleted{\unsure{fibre spéciale}} \unsure{irréduct.} Donc il en est \deleted{$X$ étant lisse sur $Y$, $X_1=X\times_Y Y_1$ l'est sur $Y_1$, $X_1$ est \unsure{normal}, donc \unsure{aussi} et} \inserted{de même de $X_1=X\times_Y Y_1$ \unsure{en $Y_1$|sur $Y_1$}, d'ailleurs le corps \unsure{résid.} \unsure{demeure}} \inserted{\unsure{celui} de $X_1$, et} le même. Il \unsure{coïncide} sur la fibre
\diagram{\begin{tikzcd}[ampersand replacement=\&,row sep=small,column sep=small]
X\arrow[dd] \& {}\\
{} \& X_1\arrow[ul]\arrow[dd]\\
Y \& {}\\
{} \& Y_1\arrow[ul]
\end{tikzcd}}
générique avec $Z_y$, il s'ensuit qu'on a un $X$-morphisme $X_1\mathrel{\deleted{\unsure{\to|\leftarrow}}}\leftarrow Z$ (qui est un isom sur les \unsure{points} de $(X_1)_y$). \deleted{\unsure{Comme} $Z$ \unsure{obtenu}, $X_1$ est \unsure{fidèle}, \unsure{puisque} sur $X$, $Y_1$ est \unsure{étale} sur $Y$.} \inserted{Par descente fid. plate, il \deleted{\unsure{faut}} s'ensuit que \deleted{\unsure{$X$}}$X_1$ est étale sur $Y$ en les pts de $(X_1)_y$, donc partout.} Donc on a
$Z\simeq X\times_Y Y_1$, $Y_1$ \unsure{revêtmt|revêtemt} fini étale de $Y$.

Par suite (revenant au cas général), \deleted{\unsure{$X_1$} est} $X_1$ est étale sur $X$ partout, et \deleted{\unsure{la comparaison}} \inserted{des considérations} de profondeur nous donnent un $X$-isom $X_1\simeq Z$, ce qui achève la démonstration.

\uline{NB} \unsure{On} a prouvé plus qu'on n'avait annoncé, savoir ceci : \unsure{On} a une suite exacte \deleted{\unsure{de}}
\margin{\unsure{TBE|TB} $\left.\rule{0pt}{3em}\right]$}
\diagram{\begin{tikzcd}[ampersand replacement=\&,column sep=small]
\pi_1(X_{\bar y},\unsure{\xi|\bar\xi})\arrow[r] \& \pi_1(X,\unsure{\xi|\zeta})\arrow[r] \& \pi_1(Y,\unsure{\zeta|\xi})\arrow[r] \& 0
\end{tikzcd}}

\unsure{Cela|Ceci} généralise (\unsure{modulo|en dehors de} l'hyp. noethérienne) SGA $\overline{\mathrm{IX}}$ 6.1., 6.2. \ldots. Ne pas confondre avec les résultats de SGA $\overline{\mathrm{X}}$, où on prend une fibre \uline{pas \inserted{\uline{néc}} générique} (en revanche, \unsure{3|?} il y faut supposer que $f$ est propre \ldots.).
% archive sheet number: 49
\scanpage{51}
\mspage{- 22 -}

n$^{\circ}$385

\[
\begin{aligned}
\langle\varphi_*(\mu),fg\rangle&=\langle\mu,(f\circ\varphi)(g\circ\varphi)\rangle=\sum\langle\mu'_i,f\circ\varphi\rangle\,\langle\mu''_i,g\circ\varphi\rangle\\
&=\sum\langle\varphi_*\mu'_i,f\rangle\,\langle\varphi_*\mu''_i,g\rangle
\end{aligned}
\]
et par suite :
\[
P\varphi_*(\mu)=\sum\varphi_*\mu'_i\otimes\varphi_*\mu''_i=(\varphi_*\otimes\varphi_*)(P(\mu)).
\]

10. \uline{Produit d'une fonction et d'une distribution}.

Soit $a$ un point de $X$. Le dual $V$ de l'algèbre commutative $C_a^r(X)$ est muni d'une structure de module sur $C_a^r(X)$ définie par la formule :
\[
\langle f.\mu,g\rangle=\langle\mu,fg\rangle\qquad\mu\in V,,f,g\in C_a^r(X).\tag*{(59)}
\]

Comme $D_a^k(X)$ est le sous-espace de $V$ formé des formes linéaires nulles sur l'\uline{idéal} $I_a^{k+1}(X)$, c'est un sous-module de $V$. Comme $D_a(X)$ est la réunion des espaces $D_a^k(X)$ pour les entiers $k\leq r$, on voit que la formule (59) définit sur $D_a(X)$ une structure de module sur l'anneau $C_a^r(X)$.

Pour tout ouvert $U$ de $X$, on définit comme suit une structure de $C^r(U)$-module sur l'ensemble des distributions de support contenu dans $U$ ; pour $f$ dans $C^r(U)$ et une distribution $\mu$ avec $S(\mu)\subset U$, on pose :
\[
(f.\mu)_a=\begin{cases}
\gamma_a(f).\mu_a & \text{si }a\in U\\
0 & \text{si }a\notin U.
\end{cases}
\]

La formule (59) reste valable pour $f,g$ dans $C^r(U)$ si $S(\mu)\subset U$.

Etant données une variété $Y$ de classe $C^r$, un morphisme $\varphi$ de $Y$ dans $X$, $\mu$ dans $D(X)$, $\unsure{\nu|\mathfrak{d}}$ dans $D(Y)$, $f$ dans $C^r(X)$ et $g$ dans $C^r(Y)$, on a les relations :
\[
(f.\mu)\times(g.\nu)=(f\times g).(\mu\times\nu)\tag*{(60)}
\]
\[
\varphi_*((f\circ\varphi).\nu)=f.\varphi_*(\nu).\tag*{(61)}
\]

PROPOSITION 6.- \uline{Soient} $X$ \uline{une variété de classe} $C^r$, $Y$ \uline{une sous-variété de} $X$, \uline{et} $\iota$
% archive sheet number: 50
\scanpage{52}
\mspage{III}

\uline{Lemme 2.} Soit $X\longrightarrow Y$ un morphisme de type fini, avec $Y$ loc. noeth. et univ. japonais. Alors pour tt faisceau d'ensembles constructible $F$ sur $X$, $f_*(F)$ est constructible.

Standart

\uline{Lemme 3} $X$ schéma connexe \deleted{\unsure{loc.}} noeth, \inserted{univ. japonais}, $(X_i)_{i\in I}$ une partition finie en parties loc fermées, $\unsure{G|C}$ un gpe fini. \deleted{Supposons que pour tout $i$, $H^1(X_i,G)$ soit fini.} \inserted{Considérons l'\deleted{application}}
\deleted{Alors \unsure{pour tt}}
\[
H^1(X,G)\longrightarrow\prod_i H^1(X_i,G)
\]

Cette application est à \uline{fibres finies}. En particulier, si pour tt $i$, $H^1(X_i,G)$ est fini, Kif Kif pour $H^1(X,G)$.

Théorie de recollement des faisceaux étales, en utilisant lemme 2. et le fait suivant :
\inserted{\uline{Lemme 4}} Si $F,G$ sont des faisceaux \inserted{d'ens} étales \unsure{constr.} sur $X$ noethérien, alors $\operatorname{Hom}(F,G)$ est un ens fini.
\margin{$\big|$}

\deleted{\uline{Proposition}} \deleted{\uline{\unsure{Corollaire|V$^{\mathrm{o}}$ Main Proposition}}} \inserted{\uline{Proposition}} $S$ \deleted{\unsure{$X$}} schéma noeth. univ. japonais, $G$ un gpe fini. Conditions équivalentes :
% archive sheet number: 51
\scanpage{53}
\mspage{- 13 -}

n$^{\circ}$385

7. \uline{Produit cartésien de deux distributions}.

Soient $X$ et $Y$ deux variétés de classe $C^r$, $a$ un point de $X$ et $b$ un point de $Y$.

Il existe une application linéaire $\iota$ de $D_a^r(X)\otimes D_b^r(Y)$ dans l'espace $B$ des formes bilinéaires sur $C_a^r(X)\times C_b^r(Y)$ caractérisée par la formule :
\[
\iota(\mu\otimes\unsure{\nu|\mathfrak{d}})(f,g)=\langle\mu,f\rangle\,\langle\nu,g\rangle.\tag*{(31)}
\]

Par ailleurs, on définit une application linéaire $\iota'$ de $D_{a,b}^r(X\times Y)$ dans $B$ par la formule :
\[
\iota'(\pi)(f,g)=\langle\pi,f\times g\rangle.\tag*{(32)}
\]

PROPOSITION 4.- \uline{Il existe une application linéaire} $\rho_{a,b}$ \uline{de} $D_{a,b}^r(X\times Y)$ \uline{dans} $D_a^r(X)\otimes D_b^r(Y)$ \uline{et une seule telle que} $\iota'=\iota\circ\rho_{a,b}$. \uline{L'application} $\rho_{a,b}$ \uline{est injective et son image est le sous-espace de} $D_a^r(X)\otimes D_b^r(Y)$ \uline{engendré par les éléments} $\mu\otimes\nu$ \uline{avec} $v(\mu)+v(\nu)\leq r$. \uline{En particulier}, \uline{si} $r$ \uline{n'est pas un entier}, \uline{l'application} $\rho_{a,b}$ \uline{est bijective}.

Nous introduirons un système de coordonnées $x_1,\ldots,x_n$ sur $X$ au voisinage de $a$ et tel que $x_i(a)=0$ pour tout $i$ ; de même, on introduit un système de coordonnées $y_1,\ldots,y_m$ sur $Y$ au voisinage de $b$ tel que $y_j(b)=0$ pour tout $j$. D'après la prop.2 il existe des bases $\{\Delta_a^\alpha\}$ de $D_a^r(X)$ et $\{\Delta_b^\beta\}$ de $D_b^r(Y)$ caractérisées par les formules :
\[
\langle\Delta_a^\alpha,x^{\alpha'}\rangle=\delta_{\alpha\alpha'}\qquad
\langle\Delta_b^\beta,y^{\beta'}\rangle=\delta_{\beta\beta'}\tag*{(33)}
\]
pour $\alpha$ et $\alpha'$ dans $N_r^n$ et $\beta$ et $\beta'$ dans $N_r^m$. Par conséquent, le produit tensoriel $D_a^r(X)\otimes D_b^r(Y)$ admet pour base les éléments $\Delta_a^\alpha\otimes\Delta_b^\beta$ pour $\alpha$ dans $N_r^n$ et $\beta$ dans $N_r^m$.

Par ailleurs (référence), on définit un système de coordonnées $z_1,\ldots,z_{n+m}$ sur $X\times Y$ au voisinage de $(a,b)$ par les conditions :
% archive sheet number: 52
\scanpage{54}
\mspage{IV}

(i) Pour tout $X$ de t. f. sur $S$, $H^1(X,\unsure{G|C})$ est fini.

(ii) Pour tt sous-schéma \deleted{\unsure{fermé}} normal \deleted{de $S$} \inserted{intègre affine} $S_0$ de $S$, \deleted{\unsure{et}} tt \unsure{rev.} fini étale $S'$ de $S_0$, tout schéma \inserted{affine} lisse de t. fini $\unsure{X|X'}$ sur $S'$, à fibres géom. connexes non vides, $H^1(\deleted{\unsure{S}}X,G)$ est fini

(iii) Pour tt $S',X$ comme dessus, $S'\ne\varnothing$, il existe un ouvert non vide $U$ de $S'$ \unsure{tel|dt|d'où} que $H^1(X\mid U',G)$ soit fini.

Il \uuline{suffit} pour ceci qu'on ait ceci :

(iv) a) Pour tout $s\in S$, et tt schéma $\bar X$ de type fini lisse affine connexe sur $\overline{k(s)}$, $H^1(\bar X,G)$ est \unsure{fini|finie}

b) Pour tout \deleted{\unsure{$S$}} \deleted{\unsure{sous}} schéma affine normal intègre $S'$ \deleted{de $S$}, \inserted{\unsure{comme|nommé} dans ii)} et tt sous \unsure{gpe|groupe} $H$ de $G\times G$, \deleted{\illegible{1}} $H^1(S',H)$ est fini.

Lorsque $G$ est commutatif, on peut affaiblir b) en y remplaçant $H$ par $G$.

\margin{\rule{4em}{.4pt}}
\uline{Dém}. Equivalence de (i) à (iii) résulte aussitôt de lemme 3. Pour prouver qu'ils sont entraînés par (iv), on doit utiliser de plus lemme 1.
% archive sheet number: 53
\scanpage{55}
\mspage{- 12 -}

n$^{\circ}$385

\uline{caractérisé par la propriété suivante} :

\uline{Pour tout} $\mu$ \uline{dans} $\bar D_a^1(X)$ \uline{et tout système de coordonnées} $x_1,\ldots,x_n$ \uline{au voisinage de} $a$, \uline{le vecteur tangent} $u(\mu)$ \uline{admet les composantes} $\langle\mu,x_1\rangle,\ldots,\langle\mu,x_n\rangle$ \uline{par rapport au système de coordonnées} $x_1,\ldots,x_n$.

Nous allons, pour terminer ce n$^{\circ}$, examiner l'effet d'un morphisme sur une distribution d'ordre 1.

Soit $\varphi$ un morp\uline{hi}sme de $X$ dans une variété $X'$ de classe $C^r$. Soit $x_1,\ldots,x_n$ un système de coordonnées au voisinage d'un point $a$ de $X$, et $x'_1,\ldots,x'_m$ un système de coordonnées sur $X'$ au voisinage du point $a'=\varphi(a)$. Le morphisme $\varphi$ s'exprime dans ces systèmes de coordonnées par les formules :
\[
x'_j(\varphi(p))=\varphi_j(x_1(p),\ldots,x_n(p))\qquad 1\leq j\leq m.\tag*{(28)}
\]

D'après les formules (15) et (24), on a alors :
\[
\langle\varphi_a\mu,x'_j\rangle=\sum_{1\leq i\leq n}\langle\mu,x_{\unsure{i|1}}\rangle\,D_i\varphi_j(x_1(a),\ldots,x_n(a))\qquad 1\leq j\leq m\tag*{(29)}
\]
pour toute distribution $\mu\in L$ ; le premier membre de cette formule représente la j-ième composante du vecteur tangent $u_{\varphi(a)}(\varphi_a(\mu))$ dans le système de coordonnées $x'_1,\ldots,x'_m$, et, par définition de l'application $T_a(\varphi)$, le second membre de (29) représente la j-ième composante du vecteur tangent $(T_a(\varphi)\circ u_a)(\mu)$ dans le même système de coordonnées sur $X'$. Autrement dit, on a :
\[
u_{\varphi(a)}\circ\varphi_a=T_a(\varphi)\circ u_a.\tag*{(30)}
\]

Nous identifierons désormais les espaces $\bar D_a^1(X)$ et $T_a(X)$ au moyen de l'isomorphisme $u_a$ précédent, lorsque $r\ne\unsure{\omega|\varpi|\infty}$.

\uline{Remarque}.- Lorsque $r=\omega$, l'espace $T_a(X)$ n'a pas été défini précédemment [et pour cause !] ; nous poserons par définition $T_a(X)=\bar D_a^1(X)$, et $T_a(\varphi)=\varphi_a$. La formule (27) reste valable si l'on définit la matrice jacobienne par $J_{ij}(a)=(D_i)_a y_j$, la base $(D_i)_a$ de $\bar D_a^1(X)$ étant définie par $\langle(D_i)_a,x_j\rangle=\delta_{ij}$.
% archive sheet number: 54
\scanpage{56}
\mspage{V}

\uline{Cor} 1 Supposons que $\unsure{G|C}$ soit premier aux car. résiduelles. \deleted{Alors \unsure{la}} Supposons que pour tt sous-schéma affine, intègre normal $S_0$ de $S$, et tt \uline{\unsure{revêtmt|revêt|ouvert}} \uline{étale} \inserted{\unsure{intègre}} $S'$ de $S_0$, enfin tt \unsure{sous-gpe|sous-groupe} $H$ de $G\times G$, on ait $H^1(S',H)$ fini. Alors pour tt $X$ de t. f. sur $S$, $H^1(X,G)$ est fini.

(\unsure{Si|Or} $G$ est commutatif, il suffit de faire $H=G$ \deleted{\unsure{et}}, et on a une condition \unsure{néc.|née} et suffisante).

\uline{Corollaire} 2 Conditions équivalentes.

(i) Pour tt $X$ de t. f. sur $Y$, et tt \unsure{grpe|gpe|groupe} \inserted{fini} $G$ premier aux car. résiduelles de $X$, on a $H^1(X\unsure{,|.}G)=\text{fini}$

(ii) \deleted{\unsure{Pour}} \unsure{Idem|Item}, avec $X$ \unsure{revêtmt|revêt} \deleted{fini} étale intègre d'un sous-schéma affine normal intègre $S_0$ de $S$.

\uline{NB} Dans le cor. 1, on a utilisé le fait que si $k$ est un corps \unsure{alg.|alg} clos, alors pour tt $X$ de t. f. sur $k$ et $G$ premier à la car. de $k$, $H^1(X,G)$ est fini. Voyons comment ceci peut être prouvé par les méthodes envisagées ici.
% archive sheet number: 55
\scanpage{57}
\mspage{- 11 -}

n$^{\circ}$385

$\Psi_i$ et $\Psi'_i$ définies et continues dans un voisinage du point $\xi=(x_1(a),\ldots,x_n(a))$ de $K^n$ et telles que :
\[
\Phi=\sum_i\Psi_i\,\Psi'_i\qquad\Psi_i(\xi)=\Psi'_i(\xi)=0
\]
et que les $\Psi_i$ soient de classe $C^r$. Ces conditions entraînent que $\Phi(\lambda_1,\ldots,\lambda_n)/\sum\limits_i|\lambda_i-\xi_i|$ tend vers $0$ si le point $(\lambda_1,\ldots,\lambda_n)$ tend vers $\xi$, et par conséquent que les dérivées de $\Phi$ sont nulles au point $\xi$. Comme on a par ailleurs $\langle\mu,f\rangle=0$ dans ce cas, la formule (24) est encore vraie.

La prop.1 montre que tout élément de $C_a^r(X)$ est congru modulo $I_a^2(X)$ au germe en $a$ d'un polynôme de degré $\leq 1$ en les $x_i$ ; les cas particuliers précédemment démontrés prouvent que (24) est vraie en général.

\hfill C.Q.F.D.

Considérons alors deux systèmes de coordonnées $x_1,\ldots,x_n$ et $y_1,\ldots,y_n$ au voisinage du point $a$ ; nous introduirons les fonctions de passage $\Phi_j$ telles que l'on ait :
\[
y_j(p)=\Phi_j(x_1(p),\ldots,x_n(p))\qquad 1\leq j\leq n\tag*{(26)}
\]
lorsque $p$ appartient à l'intersection des domaines de définition des deux systèmes de coordonnées ; par ailleurs, on introduit la matrice jacobienne par :
\[
J_{ij}(p)=D_i\,\Phi_j(x_1(p),\ldots,x_n(p)).
\]

D'après le lemme 1 appliqué à $f=y_j$, on a alors :
\[
\langle\mu,y_j\rangle=\sum_{1\leq i\leq n}\langle\mu,x_i\rangle\,J_{ij}(a)\qquad 1\leq j\leq n\tag*{(27)}
\]
pour tout $\mu\in L$. D'après la définition d'un vecteur tangent, on a finalement la proposition suivante :

PROPOSITION 3.- \uline{Soient} $a$ \uline{un point d'une variété} $X$ \uline{de classe} $C^r$ $(r\ne\unsure{\omega|\varpi|\infty})$, \uline{et} $\bar D_a^1(X)$ \uline{l'espace des distributions} $\mu$ \uline{de support} $a$ \uline{et d'ordre} $\leq 1$ \uline{telles que} $\langle\mu,1\rangle=0$. \uline{Il existe un isomorphisme} $u$ \uline{de} $\bar D_a^1(X)$ \uline{sur l'espace tangent} $T_a(X)$ \uline{à} $X$ \uline{au point} $a$,
% archive sheet number: 56
\scanpage{58}
\mspage{VI}

\deleted{On est} Récurrence sur $\dim X$, \inserted{trivial si elle est $\unsure{\leq|<}0$.} \deleted{\unsure{Utilisant}} \deleted{On \unsure{a}} Lemme 3, on est ramené au cas où $X$ est lisse et \deleted{\unsure{admet} un morphisme lisse surjectif} \inserted{connexe. Si $\dim X=1$,} c'est facile (réduction au cas complet par lemme de Abhyankar) \deleted{\unsure{cf}} c'est dû à \unsure{a} Serre-Lang. \deleted{\unsure{Dim supérieure}}

Si $\dim X\geq 2$, on est ramené \deleted{\unsure{au cas}} \inserted{par} lemme 3 au cas où il y a une fibration \inserted{lisse} $X\longrightarrow Y$, $\dim Y<\dim X$, \deleted{et} si on veut à fibres \unsure{géom.|géom} connexes non vides, de dim 1. On applique alors le lemme 1 et l'hyp. de récurrence.
\margin{?}
\uline{Prop} Soit $k$ un corps. On suppose que pour tte \deleted{\unsure{$X$}} courbe \inserted{alg.} lisse $X$ sur $k$, le \unsure{nb} des rev. étales de $X$ de degré donné soit fini. Alors la $=$ chose est vraie pour \uline{tt} $X$ de t. f. sur $k$.
% archive sheet number: 57
\scanpage{59}
\mspage{-10-}

n$^{\circ}$385

dans un voisinage ouvert de $a$. En particulier on a $\varepsilon_a=\Delta_a^{\unsure{o|0}}$.

6. \uline{Distributions et vecteurs tangents}.

La proposition 2 montre comment un système de coordonnées au voisinage d'un point $a$ de $X$ définit une base de l'espace des distributions de support $a$ ; cette base dépend du système de coordonnées choisi. Nous allons étudier comment varient les éléments d'ordre 1 de la base en fonction du système de coordonnées. \uline{Nous supposerons} $r\ne\unsure{\omega|\varpi|\infty}$.

Soient $a$ un point de $X$ et $L$ l'espace vectoriel $\bar D_a^1(X)$. Comme on a $\langle\mu,f\rangle=0$ pour toute distribution $\mu\in L$ et toute fonction constante $f$, la prop.2 montre que l'application :
\[
\mu\longrightarrow(\langle\mu,x_1\rangle,\ldots,\langle\mu,x_n\rangle)\tag*{(22)}
\]
est un isomorphisme de $L$ sur $K^n$ quel que soit le système de coordonnées $x_1,\ldots,x_n$ sur $X$ au voisinage de $a$. On peut par contre définir une base $(D_i)_a$ de $\bar D_a^1(X)$ par $\langle(D_i)_a,x_j\rangle=\delta_{ij}$. On a donc $(D_i)_a=\Delta_a^\alpha$ avec $\alpha=\varepsilon_i$. On a de plus le résultat suivant :

LEMME 1.- \uline{Soit} $f$ \uline{une fonction de classe} $C^r$ \uline{définie dans un voisinage ouvert} $V$ \uline{de} $a$ \uline{contenu dans} $U$, \uline{et soit} $\mu$ \uline{dans} $L$. \uline{Si} $f$ \uline{est de la forme} :
\[
f(p)=\Phi(x_1(p),\ldots,x_n(p))\qquad p\in V\tag*{(23)}
\]
\uline{et si} $D_i\Phi$ \uline{est la dérivée de la fonction} $\Phi$ \uline{par rapport à la} $i$-\uline{ème variable}, \uline{on a} :
\[
\langle\mu,f\rangle=\sum_{1\leq i\leq n}\langle\mu,x_i\rangle\,D_i\Phi(x_1(a),\ldots,x_n(a))\tag*{(24)}
\]
\uline{et en particulier}
\[
\langle(D_i)_a,f\rangle=D_i\Phi(x_1(a),\ldots,x_n(a))\qquad 1\leq i\leq n.\tag*{(25)}
\]

La formule (24) est évidemment satisfaite si $f=1$, ou si $f=x_i$. Supposons maintenant que le germe de $f$ en $a$ soit dans $I_a^2(X)$ ; il existe alors des fonctions
% archive sheet number: 58
\scanpage{60}
$\pi_1$ et $\underline{\operatorname{Pic}}^{\unsure{\overline{0}|\tau}}$ birationnels
% archive sheet number: 59 (second digit also resembles 3)
\scanpage{61}
\mspage{1.}

\uline{Th.} Soit $X$ préschéma alg., sur $k$ alg. clos. Alors
\[
\begin{aligned}
\pi'_1(X)={}&\text{limite projective des groupes quotients discrets de }\pi_1(X)\text{ d'\unsure{ordres|ordre}}\\
&\text{\unsure{premiers|premier} à }p
\end{aligned}
\]
est top. de type fini.

Grâce à ``van Kampen'', on peut supposer $X$ affine \inserted{\deleted{\unsure{intègre}} réduit}, puis $X$ affine et \uline{normal} [sinon, on aura
\diagram{\begin{tikzcd}[ampersand replacement=\&,column sep=small]
X'\times_X X'\arrow[r,shift left=.35ex]\arrow[r,shift right=.35ex] \& X'\arrow[r] \& X
\end{tikzcd}}
, et $\pi_1(X)$ est engendré par $\pi_1(X')$ et un nb fini de générateurs \unsure{correspondant|correspondants} aux composantes irréd. de $X'\times_X X'$, donc il en est de $=$ de $\pi'_1(X)$ \inserted{par $\pi'_1(X')$ \unsure{etc\ldots}}, au \deleted{\unsure{a}} sens des groupes \unsure{du|de} type \unsure{prodiscrets|profinis|profini}, \unsure{mot ops.|mut. ops.|mot $(p)$}

Si $\dim X=0$, trivial

Si $\dim X=1$, théorie de la spécialisation du groupe fondamental \inserted{``\unsure{tame|doux}''} nous ramène au cas de la caractéristique $0$, où on a \deleted{\unsure{un}} le \deleted{\unsure{groupe}} résultat par voie transcendante.

Si $\dim X\geq 2$, on considère une immersion
\diagram{\begin{tikzcd}[ampersand replacement=\&]
X^n\arrow[r] \& \unsure{\mathbb P|P}_k^r
\end{tikzcd}}
et un hyperplan \deleted{générique, $H$} \inserted{\unsure{variété|sous-espace} lin gén $L^{r-n+1}$}, défini sur une ext. alg close $K$ de $k$, on pose
\[
 C=\deleted{\unsure{H_{K}}}\,X^n\times_{\mathbb P_k^r}\deleted{\unsure{H}}L^{r-n+1}
\]
qui est une courbe normale de dim \unsure{un|$n$} sur $K$, et alors en vertu de Bertini ``générique'',
\diagram{\begin{tikzcd}[ampersand replacement=\&]
\pi_1(C)\arrow[r] \& \pi_1(X)
\end{tikzcd}}
est \uline{surjectif}. [\unsure{Réf|Ref}]. Donc $\pi'_1(C)\unsure{-|\longrightarrow}\pi'_1(X)$ l'est aussi, O.K.

\uline{Question} $X$ normale complète, $Y\subset X$ \deleted{\unsure{de codim} $\unsure{\geq 1|>1}$,} Alors $\pi_1(X-Y)$ mod le sous-gpe invariant fermé engendré par les groupes d'inertie wild en les composantes de \deleted{\unsure{degré}} codim 1 de $Y$, est de top. de type fini ??
% archive sheet number: 60
\scanpage{62}
\mspage{2.}

\uline{Théorème} \deleted{\unsure{Soit} $X$, $E$ \unsure{un préschéma alg. sur}} $k$ alg. clos, $X$ \deleted{\unsure{type fini sur $k$}} \deleted{complet} \inserted{\unsure{prsch.} noeth.}, \deleted{$\unsure{X'|Y}$ \unsure{le}} $D$ une partie fermée de $X$, \deleted{$\unsure{X'\to X}$} $X'$ le normalisé de $X_{\mathrm{réd}}$, $D_i$ les \deleted{composantes} \inserted{\unsure{$\xi_i$ le pt gén. de $D'_i$}} irréduct. de codim un de $p^{-1}(D)$, $U=X-D$,
\diagram{\begin{tikzcd}[ampersand replacement=\&,column sep=tiny]
 X' \arrow[d,"p"'] \& {\supset} \& \deleted{\unsure{D}} \unsure{\xi_i|D'_i}\arrow[d] \\
 X \& {\supset} \& D
\end{tikzcd}}
\deleted{$\unsure{X|Y}$} $Y$ un préschéma \unsure{noeth.} \deleted{\unsure{sur $k$}}, \deleted{\unsure{connexe} \illegible{2}} \inserted{$P=T\times U$, $P'=T\times U'$} $R$ un revêtemt étale de \deleted{$P\unsure{{}=Y\times X}$} \inserted{$P$}, \inserted{$\unsure{y|z}\in Y(k)$}, On suppose
\diagram{\begin{tikzcd}[ampersand replacement=\&,column sep=tiny]
{P'=Y\times U'} \arrow[d] \& {\subset} \& {Y\times X'=Q'} \arrow[d]\\
{\deleted{\unsure{Y\times D'_i}} \unsure{P=}Y\times U} \arrow[d] \& {\subset} \& {Y\times X=Q}\\
Y \& \&
\end{tikzcd}}

(i) $R\times_P P'$ est \inserted{\unsure{géommt}} \uline{\unsure{tamely} ramifié} \deleted{en les \unsure{fibres respectives}} \inserted{aux pts $\unsure{D_i|\xi_i}\times_k\unsure{k(y)|k(y')}$, \unsure{pt gén. des}} comp. irréduc. de $Y$. \deleted{[NB Si $Y$ est régulier, il suffit de}

(ii) $R\times_Y y$ est un rev. \uline{trivial} de $U$.
\margin{\rule{.4pt}{3em}}
Alors il existe un rev. étale $R_0$ \deleted{\unsure{sur $Y$}} \inserted{d'un voisinage de $y$ \inserted{\unsure{sur $Y'$}}} et \deleted{pour tt $y'\in Y$, $R_{y'}$ est un \unsure{rev.}} un $P$-isom $R\deleted{\unsure{(R)}}\unsure{\mid P'|\mid P}\deleted{\unsure{R_{0,Y}\times R'}}\simeq R_0\times_Y P=R_0\times U$.

\deleted{\uline{\unsure{Détail}} de $U\otimes_k k(y')$.}

On sait [réf] qu'il \unsure{suffit} de trouver pour tout $y'\in$ \deleted{\unsure{fibre infin.} \illegible{1}} \inserted{$\unsure{D_y|\mathcal O_y}\otimes\unsure{\ell|k}(\bar k)$, $R_{y'}$ est} un rev. \uline{trivial} de $k(y')$. On peut même supposer $Y$ localisé en un anneau \unsure{artin.} clos \unsure{résid.}
% slanted
\margin{Suite exacte $\pi_1^t(X,D,x)\to$ $\pi_1(X\times_k Y$}

Supposons d'abord $X$ normal, i.e. $X=X'$. \deleted{Soit} La question étant \inserted{alors} ``\unsure{birationnelle}'' sur les fibres, on peut supposer $X$ affine, même qu'on a une immersion $X\longrightarrow\mathbb P^n_k$. \deleted{\unsure{Il}}
\inserted{d'ailleurs, il suffit de regarder les $y'$ rat/$k$.} Soit $C$ une ``courbe générique'' de $X$, \deleted{\unsure{courbe} \illegible{2}} \inserted{définie} \deleted{\illegible{1}} sur une ext. alg. close $K$ de $k$. Alors
\diagram{\begin{tikzcd}[ampersand replacement=\&,row sep=tiny]
 C \arrow[r] \& X\\
 {\cup} \& {\cup}\\
 V \arrow[r] \& U
\end{tikzcd}}
$C$ est une courbe non singulière, et $\pi_1(C)\to\pi_1(X)$ est \uuline{surjectif}, d'autre part les images inverses de $D_i$ dans $C$ sont des diviseurs non singuliers, \deleted{\unsure{il en est}} \deleted{\unsure{et on vérifie} \illegible{1}}
\margin{\rule{4em}{.4pt}}
\diagram{\begin{tikzcd}[ampersand replacement=\&,row sep=tiny]
 {C\times Y=C_Y\deleted{\unsure{X}}} \arrow[r] \& Q\\
 {\cup} \& {\cup}\\
 {\unsure{U|V}\times Y=V_Y} \arrow[r] \& P
\end{tikzcd}}
On est ramené à prouver l'assertion en question en remplaçant $(k,X,U)$ par $(K,C,V)$, donc au cas d'une \uline{courbe}.

Mais alors, on est sous les conditions excellentes
% archive sheet number: 61
